\subsection{Escalabilidad}
\label{section:esca}

En el rendimiento de los programas paralelos, si aumentamos el tamaño del problema al mismo ritmo que aumentamos el número de procesadores o núcleos, y si la eficiencia no cambia (se mantiene constante), entonces nuestro programa será \emph{escalable}.
Hay un par de casos de la escalabilidad que tienen nombres especiales. Si al aumentar el número de procesadores o núcleos y podemos mantener la eficiencia constante sin aumentar el tamaño del problema, se dice que el programa es \emph{fuertemente escalable}. Por otro lado, si podemos mantener la eficiencia constante aumentando el tamaño del problema al mismo ritmo que aumentamos el número de procesadores o núcleos, entonces se dice que el programa es \emph{débilmente escalable}. \cite{pacheco2011introduction}